% ============================================================================
%  TP Final - Simulación de Sistemas (72.25) - ITBA - 2026 Q1 - Grupo G01S2
%  Nagel-Schreckenberg aplicado a vehiculos dirigidos por vibracion (Hexbug)
%
%  Compilar: pdflatex dos veces (para resolver referencias de Fig/Ec).
%  Figuras en ../figures/ (generadas por analysis/analyze.py) y fotogramas hero
%  en ../data_anim/ (generados por analysis/animate.py).
% ============================================================================
\documentclass[11pt,a4paper]{article}

\usepackage[utf8]{inputenc}
\usepackage[T1]{fontenc}
% NOTA: babel-spanish no está instalado en este TeX Live (Debian, gestionado por apt). Se omite para
% que el informe compile de forma portable; los acentos salen por inputenc/fontenc y los rótulos se
% localizan a mano abajo. En una instalación con texlive-lang-spanish puede reponerse
% \usepackage[spanish,es-tabla]{babel}.
\usepackage{lmodern}
\usepackage{microtype}
\usepackage{amsmath,amssymb}
\usepackage{graphicx}
\usepackage{float}
\usepackage{booktabs}
\usepackage[hidelinks]{hyperref}
\usepackage[margin=2.4cm]{geometry}
\usepackage{caption}
\renewcommand{\figurename}{Figura}
\renewcommand{\tablename}{Tabla}
\renewcommand{\refname}{Referencias}
\emergencystretch=1.2em % evita overfull boxes sin partición silábica española

\newcommand{\unit}[1]{\,\mathrm{#1}}
\graphicspath{{../figures/}{../data_anim/}}

\title{\textbf{Diagrama fundamental de vehículos dirigidos por vibración (VDV):
reproducción de tendencias con un autómata celular de Nagel--Schreckenberg
con interacción por contacto}}
\author{Grupo G01S2 \\
        72.25 Simulación de Sistemas --- ITBA --- 2026, primer cuatrimestre}
\date{Julio de 2026}

\begin{document}
\maketitle

\begin{abstract}
Reproducimos con un autómata celular de Nagel--Schreckenberg, en una ruta
periódica unidimensional y con la Regla 2 modificada para que la interacción sea
\emph{solo por contacto}, las tendencias del diagrama fundamental de
vehículos dirigidos por vibración (VDV, sigla técnica del artículo) medidas experimentalmente por Patterson y
Parisi~[1]. El modelo, calibrado a la geometría del experimento
($\Delta x=0{,}25\unit{mm}$, $\Delta t=1/24\unit{s}$, $\ell=44\unit{mm}$), captura
la velocidad casi constante a densidades baja y media, la caída al acercarse a la
densidad de contacto, la dependencia del orden de inserción (el orden aleatorio
queda acotado entre el creciente y el decreciente) y la convergencia de las tres
configuraciones a saturación. Discutimos con qué mecanismo el modelo alcanza (o no)
la velocidad de saturación por debajo del agente más lento, y qué rasgos del
experimento el modelo deliberadamente \emph{no} reproduce (colas de densidad por
solapamiento/desalineación). Como control, la variante clásica de la regla reproduce
el diagrama fundamental triangular analítico $Q(\rho)=\min(\rho\,v_{\max},1-\rho)$.
\end{abstract}

% ============================================================================
\section{Introducción}
% ============================================================================

El \textbf{diagrama fundamental} relaciona la velocidad media (o el flujo) de un
conjunto de partículas autopropulsadas con su densidad. Es la herramienta central
para caracterizar y diseñar sistemas de tráfico vehicular y peatonal, y en todos
ellos la velocidad decrece de forma monótona al aumentar la densidad. En el tráfico
vehicular esa caída se debe a que el conductor reduce voluntariamente su velocidad
—usando sensado remoto (la vista)— para mantener distancia; en peatones, en cambio,
a densidades altas aparece una segunda causa de naturaleza \emph{física}: el contacto
y el empuje entre cuerpos. Aislar experimentalmente esa contribución puramente de
contacto es difícil porque en general convive con el sensado remoto.

Patterson y Parisi~[1] estudian un sistema mínimo que interactúa \emph{exclusivamente
por contacto}, sin sensado remoto: robots comerciales Hexbug Nano (VDV, cuerpo de
$44\times15\times18\unit{mm}$) confinados en un canal circular de
$18\unit{mm}$ de ancho y $\approx1313\unit{mm}$ de longitud efectiva, que admite hasta
$30$ vehículos. Cada robot avanza por la rectificación de su vibración sobre cerdas
asimétricas; caracterizados individualmente, sus velocidades libres se distribuyen de
forma aproximadamente uniforme en $[90,120]\unit{mm/s}$. El experimento se filma a
$24\unit{fps}$ y se mide la velocidad de cada agente por diferencias finitas.

Sus hallazgos, que este trabajo se propone reproducir como \emph{tendencias}, son:
(i) la velocidad media permanece casi constante a densidad baja y media y cae entre
un $25\%$ y un $40\%$ al acercarse la densidad lineal a la inversa de la distancia de
contacto; (ii) el resultado depende del \emph{orden} en que se introducen los agentes
—ordenados por su velocidad libre— y el orden aleatorio queda acotado entre el
creciente y el decreciente; (iii) a densidad máxima ($N=30$) las tres configuraciones
convergen a una velocidad similar; y (iv) esa velocidad de saturación es \emph{menor}
que la del agente más lento ($90\unit{mm/s}$), lo que evidencia un colapso colectivo.

El modelo general de Nagel--Schreckenberg~[2] describe el tráfico como un autómata
celular estocástico: los vehículos ocupan celdas de una pista, con velocidad entera
acotada, y se actualizan de forma sincrónica según reglas locales de aceleración,
frenado aleatorio, interacción y movimiento. Adoptamos ese marco por su bajo costo y
su capacidad conocida de generar diagramas fundamentales, y modificamos su regla de
interacción para que sea de \emph{contacto puro}, acorde a la física del experimento.

% ============================================================================
\section{Modelo}
% ============================================================================

La pista es una \textbf{ruta unidimensional con condiciones periódicas} de $L$
celdas: lo que sale por un extremo reentra por el otro. El canal circular del
experimento se modela así como una línea periódica; la curvatura no interviene en el
modelo. Cada vehículo $i$ es una partícula \emph{extendida} que ocupa $\ell$ celdas,
con posición de cola $x_i$ y velocidad entera $v_i\in\{0,\dots,v_{\max,i}\}$
(celdas/paso). El hueco (celdas libres) hasta el vehículo de adelante es

\begin{equation}
g_i = \left(x_{i+1} - x_i - \ell\right) \bmod L .
\label{eq:gap}
\end{equation}

La actualización es \textbf{sincrónica} y se aplica en el orden
$R1 \to R3 \to R2 \to R4$. Aplicar el frenado aleatorio (R3) \emph{antes} de proyectar
los contactos (R2) es deliberado: el frenado solo puede reducir el avance, de modo que
proyectar después garantiza que nunca aparezca un solapamiento que el frenado no
pudiera prever. Las reglas son

\begin{align}
\text{(R1) Aceleración:}\quad & v_i \leftarrow \min(v_i + 1,\, v_{\max,i}) \label{eq:r1}\\
\text{(R3) Frenado aleatorio:}\quad & v_i \leftarrow \max(0,\, v_i - 1)\ \text{con probabilidad } p \label{eq:r3}\\
\text{(R2) Colisión:}\quad & \text{Ec.~\eqref{eq:r2a} (contacto puro) o \eqref{eq:r2b} (clásica)} \nonumber\\
\text{(R4) Movimiento:}\quad & x_i \leftarrow (x_i + d_i) \bmod L \label{eq:r4}
\end{align}

donde $d_i$ es el desplazamiento efectivo del paso, definido por R2. Con $p=0$ la regla
R3 nunca se aplica y el orden se reduce a $R1\to R2\to R4$; solo la variante clásica
(B) coincide con el Nagel--Schreckenberg canónico y su diagrama triangular.

En la variante oficial \textbf{contacto puro} (A) el vehículo avanza libremente hasta
alcanzar al de adelante; si lo alcanzaría, queda inmediatamente detrás (a contacto,
$g=0$, sin solaparse) y \emph{hereda} su velocidad. Sea $\hat v_i$ la velocidad deseada
tras R1 y R3; los desplazamientos se resuelven, de adelante hacia atrás, como el mayor
punto fijo

\begin{equation}
d_i = \min\!\left(\hat v_i,\; g_i + d_{i+1}\right),
\label{eq:r2a}
\end{equation}

con $d_{i+1}$ el desplazamiento del líder en el mismo paso. Esta relajación por
agrupamientos garantiza el \textbf{no-solapamiento}: todo desplazamiento cumple
$d_i \le g_i + d_{i+1}$, de modo que el hueco resultante nunca es negativo. La velocidad
que se hereda para el paso siguiente es la del líder cuando se llega a contacto, y $d_i$
en caso contrario.

La variante \textbf{clásica} (B), usada solo para validar el motor
(Sec.~\ref{sec:validacion}), es la de Nagel--Schreckenberg de libro: se anticipa la
posición del líder,
\begin{equation}
v_i \leftarrow \min(\hat v_i,\, g_i)\,;\qquad g_i=0 \Rightarrow v_i \leftarrow \min(\hat v_i, d_{i+1}).
\label{eq:r2b}
\end{equation}

Cada vehículo tiene una \textbf{velocidad libre} propia
$v^{\mathrm{libre}}_i \sim \mathcal{U}[90,120]\unit{mm/s}$, heterogénea como en el
experimento, que fija su velocidad máxima entera $v_{\max,i}$ (Sec.~\ref{sec:simulaciones}).

% ============================================================================
\section{Implementación}
% ============================================================================

El motor está implementado en Java. Mantiene los vehículos en una lista ordenada
cíclicamente sobre la ruta y ejecuta cada paso sincrónico sobre una \emph{instantánea
inmutable} de huecos y velocidades deseadas, de modo que el resultado no depende del
orden de iteración. La regla de colisión recibe esa instantánea y devuelve, por vehículo,
un par $(\text{desplazamiento } d_i,\ \text{velocidad heredada})$; no muta el estado.
El paso, en pseudocódigo:

\begin{verbatim}
paso():
    si protocolo incremental y toca lote: insertar 5 vehículos (por empuje)
    para cada vehículo i:  deseada_i <- min(v_heredada_i + 1, vmax_i)     # R1
                           con prob p: deseada_i <- max(0, deseada_i - 1)  # R3
    movimientos <- ReglaColisión.resolver(instantanea(huecos, deseadas))  # R2
    para cada vehículo i:  x_i <- (x_i + d_i) mod L                        # R4
                           registrar d_i como velocidad de cuadro
                           registrar velocidad heredada para el próximo paso
\end{verbatim}

El motor distingue dos velocidades por vehículo: la \emph{velocidad de cuadro} $d_i$
(el desplazamiento efectivo, lo que mediría una cámara por diferencia de posiciones,
que es la que se emite) y la \emph{velocidad heredada} que alimenta R1 del paso
siguiente. En la variante clásica ambas coinciden; en contacto puro difieren solo en
el cuadro en que un vehículo queda a contacto (avanza hasta el líder pero hereda su
velocidad). El motor es 100\% determinista dado el identificador de la realización:
toda la aleatoriedad proviene de un único generador inicializado con ese identificador,
y con $p=0$ no se consume ningún número aleatorio en la dinámica.

Para el protocolo incremental, agregar un lote de vehículos se hace por
\textbf{empuje}: el vehículo nuevo (velocidad inicial $0$) se coloca a contacto delante
de uno existente y los que le siguen se corren hacia adelante lo mínimo necesario para
no solaparse, consolidando el espacio libre de la ruta. Esto permite alcanzar $N=30$,
donde $30\,\ell = L$ y la ruta queda exactamente llena a contacto. El constructor de la
ruta valida el no-solapamiento al armar la condición inicial y en cada inserción; en los
pasos ordinarios el no-solapamiento queda garantizado por construcción de la Regla 2
(Ec.~\eqref{eq:r2a}, $d_i \le g_i + d_{i+1}$). La CLI registra el estado
\emph{antes} de ejecutar cada paso, por lo que un lote insertado en el paso $k$ se observa
a partir de la muestra $k+1$; las fases del protocolo incremental quedan definidas por el
$N$ activo real de cada cuadro.

Siguiendo el criterio de la cátedra, el motor emite \textbf{solo variables físicas}
(identificador de paso, id de vehículo, $x$ en mm y $v$ en mm/s). No guarda color ni
radio: el color de las animaciones se deriva de la velocidad \emph{después} de la
simulación. Todos los observables se calculan por fuera del motor.

% ============================================================================
\section{Simulaciones}\label{sec:simulaciones}
% ============================================================================

\subsection{Calibración y geometría}

La malla discreta se calibra a la geometría del experimento
(Tabla~\ref{tab:calib}). El cuanto de velocidad es $\Delta v=\Delta x/\Delta t=6\unit{mm/s}$
y la velocidad máxima de cada vehículo es $v_{\max,i}=\mathrm{round}(v^{\mathrm{libre}}_i/\Delta v)\in\{15,\dots,20\}$.
La longitud de pista se fija en $L=30\,\ell=5280$ celdas $=1320\unit{mm}$ para que
$N=30$ cierre exactamente a contacto (el experimento reporta $\approx1313\unit{mm}$;
la diferencia se discute en Sec.~\ref{sec:limitaciones}). La densidad global es
$\rho=N/L_{\mathrm{fis}}$, que a $N=30$ vale $0{,}0227\unit{mm^{-1}}=1/(44\unit{mm})$.

\begin{table}[h]\centering
\caption{Parámetros de calibración (fijos en todas las corridas).}
\label{tab:calib}
\begin{tabular}{lll}
\toprule
Cantidad & Símbolo & Valor \\
\midrule
Paso espacial & $\Delta x$ & $0{,}25\unit{mm}$ \\
Paso temporal & $\Delta t$ & $1/24\unit{s}$ \\
Cuanto de velocidad & $\Delta v=\Delta x/\Delta t$ & $6\unit{mm/s}$ \\
Largo de vehículo & $\ell$ & $176$ celdas $=44\unit{mm}$ \\
Largo de pista & $L$ & $5280$ celdas $=1320\unit{mm}$ \\
Velocidad máxima & $v_{\max,i}=\mathrm{round}(v^{\mathrm{libre}}_i/\Delta v)$ & $\{15,\dots,20\}$ \\
\bottomrule
\end{tabular}
\end{table}

\subsection{Parámetros, condiciones iniciales y protocolos}

Los \textbf{parámetros} que se varían son el número de vehículos
$N\in\{5,10,15,20,25,30\}$ y la probabilidad de frenado aleatorio
$p\in\{0,\,0{,}1,\,0{,}2,\,0{,}3,\,0{,}4\}$. Las \textbf{condiciones iniciales} de
cada realización —distintas de los parámetros— son las posiciones iniciales (sin
solaparse), el valor concreto de cada velocidad libre y la secuencia reproducible del
generador; todos los vehículos arrancan con velocidad $0$. Una \emph{realización} queda
determinada por su identificador reproducible.

Se usan dos protocolos:
\begin{itemize}
  \item \textbf{$N$ fijo}: corridas independientes con $N$ constante, para barrer
  $N$ y $p$ de forma limpia. $10\,000$ pasos ($\approx417\unit{s}$).
  \item \textbf{Incremental cada $180\unit{s}$}: réplica del protocolo del artículo.
  Se arranca con $5$ vehículos y se agregan $5$ cada $180\unit{s}$ hasta $N=30$, con
  $25\,920$ pasos $=6\times180\unit{s}=1080\unit{s}$. Se examinan tres
  \textbf{órdenes de inserción} según la velocidad libre: creciente
  (los más lentos primero), decreciente (los más rápidos primero) y aleatorio.
\end{itemize}
Se corren $M=30$ \textbf{realizaciones} por punto. El estado se guarda cada $10$ pasos.

\subsection{Observables}

Todos se calculan \textbf{después} de la simulación, a partir del estado
$(x_i,v_i)$ guardado. La velocidad media de una realización es el promedio sobre los
$N$ vehículos y los $T$ pasos del régimen estacionario,
\begin{equation}
\bar{v} = \frac{1}{N\,T}\sum_{t}\sum_{i=1}^{N} v_i(t) \quad [\unit{mm/s}],
\label{eq:vmedia}
\end{equation}
equivalente al observable global $\langle L_i/T\rangle$ del artículo (la velocidad de
cuadro es el desplazamiento real por paso, sin ambigüedad de envoltura). Se informa con
su \textbf{desvío entre las $M$ realizaciones} —no el de todos los cuadros juntos, que
subrepresenta el error—:
\begin{equation}
\sigma_{\bar v} = \sqrt{\frac{1}{M-1}\sum_{m=1}^{M}\left(\bar v^{(m)} - \langle \bar v\rangle\right)^2}.
\label{eq:error}
\end{equation}
La densidad individual es la inversa de la distancia (centro a centro) al vecino más
cercano sobre la ruta,
\begin{equation}
\rho_i = \frac{1}{\Delta x\,(\ell + g_i^{\,\mathrm{vecino}})} \quad [\unit{mm^{-1}}],
\label{eq:densidad}
\end{equation}
con pico esperado en $1/(44\unit{mm})\approx0{,}0227\unit{mm^{-1}}$ (contacto). Se reportan
las distribuciones (PDF) de $\rho_i$ y de la velocidad microscópica, y el diagrama
fundamental (velocidad instantánea vs $\rho_i$, ordenado por densidad y suavizado con
media móvil). El \textbf{corte del estacionario} es un único valor global elegido
\emph{por inspección} de la evolución temporal (no un descarte fijo en porcentaje),
aplicado de forma uniforme a todos los $(N,p)$ y registrado en
\texttt{figures/manifiesto.csv}. En el protocolo incremental, por su diseño, el
observable se promedia sobre la ventana completa de $180\unit{s}$ de cada fase
($\langle L_i/180\unit{s}\rangle$, igual que el artículo).

% ============================================================================
\section{Resultados}
% ============================================================================

\subsection{Estudio con $N$ fijo}

La Figura~\ref{fig:anim-fixed} muestra fotogramas de la dinámica en régimen a baja
densidad ($N=5$, flujo libre) y a saturación ($N=30$, ruta llena a contacto). A baja
densidad los vehículos avanzan casi sin interactuar; a $N=30$ forman un anillo rígido a
contacto.

\begin{figure}[H]\centering
\includegraphics[width=0.9\linewidth]{N5_p01_CONTACTO_PURO_FIXED_N_r1_fotograma.png}\\[4pt]
\includegraphics[width=0.9\linewidth]{N30_p01_CONTACTO_PURO_FIXED_N_r1_fotograma.png}
\caption{Fotogramas en régimen a baja densidad ($N=5$, arriba) y a saturación
($N=30$, abajo), contacto puro, $p=0{,}1$. Color según velocidad. Se muestra un fotograma
representativo; la animación completa se regenera con \texttt{scripts/generar\_entrega.sh}.}
\label{fig:anim-fixed}
\end{figure}

El escalar que caracteriza cada corrida es el promedio estacionario de la velocidad
media (Ec.~\ref{eq:vmedia}). La Figura~\ref{fig:evol-fixed} muestra la evolución
temporal de la velocidad media y el corte del transitorio elegido por inspección; a
partir de ese corte la serie fluctúa alrededor de un valor estable.

\begin{figure}[H]\centering
\includegraphics[width=0.7\linewidth]{evolucion_temporal_contacto_puro_sin_orden_fixed.png}
\caption{Evolución temporal de la velocidad media ($N=30$, contacto puro). La línea
punteada marca el corte del transitorio elegido por inspección; el observable reportado
es el promedio a partir de ese corte (Ec.~\ref{eq:vmedia}).}
\label{fig:evol-fixed}
\end{figure}

La curva respuesta--estímulo (Figura~\ref{fig:vN-fixed}) resume la velocidad media en
función de $N$, una curva por $p$. A densidad baja la velocidad media se mantiene
cercana a la velocidad del agente más lento presente ($\approx90$--$95\unit{mm/s}$, un
$\sim10\%$ por debajo de la media libre poblacional de $105\unit{mm/s}$) y depende poco
de $N$: en la pista periódica de un solo carril, sin adelantamiento y con contacto puro,
a densidad baja los vehículos corren mayormente libres pero son alcanzados y retenidos una y
otra vez por el más lento —el agrupamiento a contacto no es permanente—, de modo que la
velocidad \emph{media} queda anclada a la del mínimo de los $N$ sorteos de $v_{\max}$. La caída suave de
$\bar v(N)$ proviene de ese agrupamiento: al crecer $N$ es más probable muestrear un
agente aún más lento y el cúmulo lo sigue; es el mismo mecanismo que satura en $N=30$. Al
acercarse a la ruta llena la velocidad cae. En el punto singular $N=30$ (ruta exactamente
llena) el valor de saturación depende de $p$: con $p=0$ el anillo rígido cruza a la
velocidad del más lento, y con $p>0$ el frenado aleatorio del cúmulo lo hace caer por
debajo de ella. Ese frenado es abrupto en el anillo lleno exacto: a $p=0{,}2$ satura en
$\approx45\pm14\unit{mm/s}$ (dispersión entre realizaciones), pero a $p\ge0{,}3$ el sistema
entra en un régimen \textbf{congelado} ($v\approx0$): la mayoría de las realizaciones quedan
detenidas y la mediana es $\approx0\unit{mm/s}$ (la media puntual, $\approx0{,}5\unit{mm/s}$ a
$p=0{,}3$, no es un valor típico), por el ``trinquete'' de que los $30$ vehículos deben no
frenar en el mismo paso para avanzar ($(1-p)^{30}\approx0$). Ese congelamiento no corresponde a ningún régimen del experimento
—que satura en $\approx70\unit{mm/s}$ y nunca se congela—; el punto físicamente comparable
es $p=0{,}1$.

\begin{figure}[H]\centering
\includegraphics[width=0.7\linewidth]{velocidad_media_vs_N_contacto_puro_sin_orden_fixed.png}
\caption{Velocidad media en función de $N$, una curva por probabilidad de frenado $p$
(contacto puro, $N$ fijo). $N=30$ es la ruta llena exacta (punto singular); allí el
descenso de la velocidad depende de $p$.}
\label{fig:vN-fixed}
\end{figure}

Las distribuciones microscópicas (Figura~\ref{fig:pdf-fixed}) muestran el pico de
densidad individual en $1/(44\unit{mm})$ (geometría de contacto) acentuándose al aumentar
$N$, y la PDF de velocidad angostándose y corriéndose a valores menores a alta densidad.

\begin{figure}[H]\centering
\includegraphics[width=0.49\linewidth]{pdf_densidad_p0.1_contacto_puro_sin_orden_fixed.png}\hfill
\includegraphics[width=0.49\linewidth]{pdf_velocidad_p0.1_contacto_puro_sin_orden_fixed.png}
\caption{Distribuciones (PDF) de densidad individual (izquierda; pico en $1/(44\unit{mm})$)
y de velocidad microscópica (derecha), una curva por $N$ ($N$ fijo, contacto puro,
$p=0{,}1$).}
\label{fig:pdf-fixed}
\end{figure}

El diagrama fundamental con $N$ fijo (Figura~\ref{fig:fd-fixed}) resume velocidad vs
densidad individual: velocidad casi plana hasta cerca de la densidad de contacto y luego
descenso.

\begin{figure}[H]\centering
\includegraphics[width=0.7\linewidth]{diagrama_fundamental_contacto_puro_fixed_n.png}
\caption{Diagrama fundamental (velocidad vs densidad individual $\rho_i$) con $N$ fijo,
una curva por $p$ (contacto puro).}
\label{fig:fd-fixed}
\end{figure}

\subsection{Estudio incremental: efecto del orden de inserción}

Este protocolo reproduce el del artículo. La Figura~\ref{fig:anim-inc} muestra
fotogramas de los tres órdenes de inserción.

\begin{figure}[H]\centering
\includegraphics[width=0.9\linewidth]{INC_p01_CONTACTO_PURO_ASCENDING_r1_fotograma.png}\\[3pt]
\includegraphics[width=0.9\linewidth]{INC_p01_CONTACTO_PURO_DESCENDING_r1_fotograma.png}\\[3pt]
\includegraphics[width=0.9\linewidth]{INC_p01_CONTACTO_PURO_RANDOM_r1_fotograma.png}
\caption{Fotogramas del protocolo incremental para los órdenes creciente (arriba),
decreciente (medio) y aleatorio (abajo) (contacto puro, $p=0{,}1$). Se muestra un fotograma
representativo; la animación completa se regenera con \texttt{scripts/generar\_entrega.sh}.}
\label{fig:anim-inc}
\end{figure}

La evolución temporal (Figura~\ref{fig:evol-inc}) muestra la estructura en escalera: cada
$180\unit{s}$ se insertan $5$ vehículos y la velocidad media da un salto y se reacomoda.
Cada fase de $180\unit{s}$ define un valor de $N$ activo; el observable se promedia sobre
la ventana completa de cada fase, como en el artículo.

\begin{figure}[H]\centering
\includegraphics[width=0.7\linewidth]{evolucion_temporal_incremental_contacto_puro_p01_incremental.png}
\caption{Evolución temporal de la velocidad media en el protocolo incremental, una curva
por orden de inserción (contacto puro, $p=0{,}1$). Los escalones marcan las inserciones
cada $180\unit{s}$.}
\label{fig:evol-inc}
\end{figure}

La curva respuesta--estímulo por orden (Figura~\ref{fig:vN-inc}, análoga a la Fig.~2 del
artículo) muestra que a densidad baja y media el orden decreciente da mayor velocidad que
el creciente, con el aleatorio en el medio, y que las tres tienden a acercarse a
saturación.

\begin{figure}[H]\centering
\includegraphics[width=0.7\linewidth]{velocidad_media_incremental_contacto_puro_p01_incremental.png}
\caption{Velocidad media en función de $N$ para los tres órdenes de inserción (contacto
puro, $p=0{,}1$; protocolo incremental). Comparar con la Fig.~2 del artículo~[1].}
\label{fig:vN-inc}
\end{figure}

Las distribuciones por orden (Figura~\ref{fig:pdf-inc}, análogas a las Figs.~3 y 4 del
artículo) separan las fases por $N$ activo: el pico de densidad se corre hacia el contacto
al aumentar $N$. En el orden creciente (el mostrado) la PDF de velocidad es casi
\emph{independiente} de $N$: el agente más lento está presente desde $N=5$, de modo que la
distribución queda anclada en velocidades bajas ($\approx90\unit{mm/s}$). El corrimiento
marcado con $N$ aparece en cambio en el orden decreciente, que parte de velocidades altas
($\approx116\unit{mm/s}$ a $N=5$) y se desplaza hacia abajo al ir agregándose los agentes
más lentos (hasta $\approx81\unit{mm/s}$ a $N=30$).

\begin{figure}[H]\centering
\includegraphics[width=0.49\linewidth]{pdf_densidad_incremental_contacto_puro_ascending_p01_incremental.png}\hfill
\includegraphics[width=0.49\linewidth]{pdf_velocidad_incremental_contacto_puro_ascending_p01_incremental.png}
\caption{Distribuciones por $N$ activo en el protocolo incremental, orden creciente
(contacto puro, $p=0{,}1$): densidad individual (izquierda) y velocidad microscópica
(derecha). Análogas a las Figs.~3A y 4A del artículo~[1]. Las figuras de los órdenes
decreciente y aleatorio se generan de igual forma.}
\label{fig:pdf-inc}
\end{figure}

Finalmente, el diagrama fundamental por orden (Figura~\ref{fig:fd-inc}, análogo a la
Fig.~5D del artículo) confirma la estructura: el orden aleatorio queda acotado entre el
creciente y el decreciente.

\begin{figure}[H]\centering
\includegraphics[width=0.7\linewidth]{diagrama_fundamental_contacto_puro_incremental_180s_p0.1.png}
\caption{Diagrama fundamental por orden de inserción (contacto puro, $p=0{,}1$;
protocolo incremental). El orden aleatorio queda acotado entre el creciente y el
decreciente. Comparar con la Fig.~5D del artículo~[1].}
\label{fig:fd-inc}
\end{figure}

% ============================================================================
\section{Validación}\label{sec:validacion}
% ============================================================================

Como control del motor usamos la variante clásica (B, Ec.~\ref{eq:r2b}) en su régimen de
validez: partículas puntuales ($\ell=1$), velocidad máxima homogénea, frenado $p=0$ y
condición inicial de reparto uniforme. En ese caso el estado estacionario es
$\bar v = \min(v_{\max}, g)$ y el flujo $Q=\rho\,\bar v$ debe coincidir con el diagrama
fundamental triangular analítico
\begin{equation}
Q(\rho)=\min(\rho\,v_{\max},\,1-\rho).
\label{eq:triangular}
\end{equation}
El motor reproduce esta curva a precisión de máquina (Figura~\ref{fig:validacion};
verificado además en los tests automatizados con tolerancia $10^{-9}$, cubriendo flujo
libre, el pico $\rho_c=1/(v_{\max}+1)$ y la rama congestionada). La variante de
\emph{contacto puro} \textbf{no} se valida contra esta curva: sin anticipación, mantiene
flujo libre ($Q=\rho\,v_{\max}$) hasta el contacto, por lo que su diagrama es distinto del
triangular por construcción.

\begin{figure}[H]\centering
\includegraphics[width=0.62\linewidth]{validacion_triangular.png}
\caption{Validación de la variante clásica contra el diagrama fundamental triangular
analítico $Q(\rho)=\min(\rho\,v_{\max},1-\rho)$ ($v_{\max}=5$, $p=0$, $\ell=1$, reparto
uniforme). Los puntos simulados caen sobre la curva analítica (error máximo
$\sim10^{-16}$).}
\label{fig:validacion}
\end{figure}

% ============================================================================
\section{Limitaciones}\label{sec:limitaciones}
% ============================================================================

El modelo reproduce \emph{tendencias}, no el experimento punto a punto. Sus límites,
que enmarcan las conclusiones:
\begin{itemize}
  \item \textbf{Sin colas de densidad más allá del contacto.} El experimento observa
  densidades $\rho_i>1/(44\unit{mm})$ por solapamiento parcial, desalineación por la
  geometría triangular del VDV y errores de seguimiento. El modelo impone
  no-solapamiento duro, así que \emph{deliberadamente} no reproduce esa cola: su PDF de
  densidad se corta en $1/(44\unit{mm})$.
  \item \textbf{Punto $N=30$ singular y dependiente de $L$.} Se fija $L=1320\unit{mm}$
  ($=30\ell$) para que $N=30$ cierre exacto a contacto; el experimento tiene
  $\approx1313\unit{mm}$, físicamente más corto que $30$ vehículos rígidos, lo que en la
  realidad admite solapamiento pero en el modelo es imposible. La sensibilidad a $L$ solo
  puede correrse para $N\le29$ (a $L=1313$, $30\ell=5280>5252$ celdas no entran): el punto
  de saturación $N=30$ solo existe a $L=1320$.
  \item \textbf{Mecanismo del colapso a saturación.} La coincidencia con ``velocidad de
  saturación menor que la del más lento'' se apoya en el punto lleno-exacto $N=30$ y
  requiere $p>0$ (el frenado aleatorio del anillo rígido). No es el mismo mecanismo físico
  del experimento (agrupamiento con solapamiento/desalineación): es el frenado estocástico
  de un cúmulo rígido en el cierre exacto de la ruta. Se reporta como reproducción de la
  \emph{tendencia}, no del mecanismo.
  \item \textbf{Velocidad cuantizada.} La velocidad está discretizada en pasos de
  $\Delta v=6\unit{mm/s}$; el modelo no reproduce colas instantáneas finas de velocidad.
  \item \textbf{Forma de la PDF de velocidad.} El modelo reproduce la tendencia de
  angostamiento y corrimiento con $N$, pero no la forma: su PDF son deltas cuantizadas
  ($\Delta v=6\unit{mm/s}$) con la moda entre $\approx84$ y $90\unit{mm/s}$ (la velocidad
  del cúmulo bloqueado por el más lento, en saturación un cuanto por debajo de su velocidad
  libre), mientras el experimento tiene una distribución suave con pico en
  $\approx65\unit{mm/s}$ y cola hasta $\approx300\unit{mm/s}$. Al agruparse rígidamente
  detrás del más lento, la velocidad microscópica del modelo queda concentrada cerca de la
  del cúmulo y \emph{no} reproduce la población ancha de baja velocidad
  ($\approx30$--$65\unit{mm/s}$) que muestra el experimento.
  \item \textbf{$p$ representativo.} Las figuras microscópicas se muestran a $p=0{,}1$; no
  son automáticamente representativas de todo $p$ (el barrido completo en $p$ está en los
  datos).
\end{itemize}

% ============================================================================
\section{Conclusiones}
% ============================================================================

Sobre lo mostrado en Resultados:
\begin{itemize}
  \item El autómata celular con interacción por contacto reproduce las tendencias
  centrales del diagrama fundamental de VDV: velocidad casi constante a densidad baja y
  media y caída al acercarse a la densidad de contacto
  (Figs.~\ref{fig:fd-fixed},~\ref{fig:fd-inc}).
  \item El orden de inserción modula el diagrama en el protocolo incremental: el orden
  aleatorio queda acotado entre el creciente y el decreciente, y las tres configuraciones
  se acercan a saturación (Figs.~\ref{fig:vN-inc},~\ref{fig:fd-inc}), como en el
  artículo~[1].
  \item El pico de densidad individual en $1/(44\unit{mm})$ y el angostamiento de la PDF de
  velocidad con $N$ reproducen la \emph{tendencia} de angostamiento y corrimiento con $N$,
  con la salvedad de que el modelo no genera la cola de densidad por solapamiento ni la
  forma detallada de la PDF (Sec.~\ref{sec:limitaciones}).
  \item La caída de la velocidad de saturación por debajo del agente más lento se obtiene
  en el punto singular $N=30$ y requiere frenado aleatorio ($p>0$); es la reproducción de
  una tendencia, con un mecanismo distinto al experimental (Sec.~\ref{sec:limitaciones}).
  \item La variante clásica valida el motor contra el diagrama fundamental triangular
  analítico con tolerancia $10^{-9}$ (Sec.~\ref{sec:validacion}).
\end{itemize}

\clearpage
% ============================================================================
\section*{Referencias}
\addcontentsline{toc}{section}{Referencias}
% ============================================================================
\begin{enumerate}
\item G.~A.~Patterson, D.~R.~Parisi, ``Fundamental diagram of vibration-driven
vehicles'', preprint, 31 de octubre de 2023.
\item K.~Nagel, M.~Schreckenberg, ``A cellular automaton model for freeway traffic'',
\emph{Journal de Physique I}, 2(12), 2221--2229 (1992).
\end{enumerate}

\end{document}
